\def\ChapterCode{EIN}
\chapter
[\CO{[EIN]} Einführung in die Medizin]
{Einführung in die Medizin}
\label{chp:EIN}

\ChapterIntro
{%
Die Medizin ist die Lehre von Gesundheit und
Krankheit sowie die Kunde der Gesunderhaltung, Heilung und
Linderung.
}
{%
\begin{description}
  \item[Maintainer:] --
  \item[Autoren:] Diverse
  \item[Reviewer:] Standard-Reviewprozess
  \item[Version:] {\VarDocVersionTypeName} ({\VarDocVersionTypeDescription})
  \item[SHA1:] {\DocGitCommit}
\end{description}

{\TxTeilkapitelAASS}
}



\clearpage % TEMP temp clearpage
\section{\HeadingKeyword{Begriffe}}

\nocite{BrockhausDrei1Bd2}
\nocite{BrockhausDrei2Bd2}
\nocite{BrockhausDrei3Bd2}
\nocite{BertelsmannVolkslexikon25}
\nocite{Pschy259}

\Dictum
[Anton Kuh]
{Alles Unglück kommt von der Terminologie.}

Die Kenntnis von Fachbegriffen dient der Kommunikation und
Verständigung mit anderem Fachpersonal.


\subsection{\HeadingKeyword{Allgemeine} Begriffe}


% \Text{Symbole}{
% 
% % \begin{multicols}{2}
% \begin{description}
% \end{description}
% % \end{multicols}
% }{}{}{}{}


\Text{Klinik -- Diagnose -- Therapie}
{

%\begin{multicols}{2}
\begin{description}
  \item[\T{Symptom}]
  \index{Symptom}
  Krankheitserscheinung
  
  \item[\T{Syndrom}]
  \index{Syndrom}
  definierte Kombination von Symptomen
  
  \item[\T{Komplikationen}]
  \index{Komplikation}
  möglicherweise auftretende Schwierigkeiten
  \item[\T{Therapie}] \index{Therapie} Ma{\ss}nahmen zur
  Heilung oder Linderung von Krankheiten und Symptomen
  
  \item[\T{Indikation}] Grund zur Anwendung eines bestimmten
  Verfahrens.
  \item[Kontraindikation] Grund zur \E{Nicht}\-an\-wendung
  eines bestimmten Verfahrens. Es gibt \E{relative} und \E{absolute}
  Kontraindikationen. Relative Kontraindikationen treffen nur
  unter bestimmten Umständen oder in gewissem Umfang zu,
  absolute Kontraindikationen verbieten das Verfahren
  praktisch immer.
  \item[\T{Diagnose}] Zuordnung von Beschwerden und Symptomen zu
  einem definierten Krankheitsbild. Es gibt verschiedene Arten
  und Funktionen von Diagnosen, siehe
  \FullPrettyRef{topic:diagnose}.
  \item[\T{Diagnostik}]\index{Diagnostik}
  
  Als Diagnostik versteht man alle Bemühungen um eine
  Diagnose zu stellen, egal ob man sich der eigenen Sinne
  oder komplizierter Geräte bedient.
  
  \item[\T{Befund}]\index{Befund}
  
  Jede diagnostische Maßnahme hat (hoffentlich) ein
  Ergebnis, welches man als
  \T{Befund} bezeichnet.
  
  \item[\T{Morbus}] \index{Morbus}\index{M.}
  \Lang{Abk.} M. \Lang{lat.}: Krankheit, Erkrankung.
  %
  Die Bezeichnung Morbus wird oft zusammen mit einem Eigenamen, 
  zumeist vom Entdecker einer Erkankung, 
  oder einem lateinische Zusatz 
  zur Benennung von Krankheiten verwendet; z.\,B.: 
  Morbus Alzheimer (Demenz vom Alzheimer-Typ), 
  Morbus Koch (Tuberkulose), 
  Morbus Parkinson, 
  Morbus comitialis (Epilepsie).
\end{description}
%\end{multicols}
}
{}
{}
{}
{}




\Text{Vorsilben}{

\begin{multicols}{2}
\begin{description}
  \item[Hyper-] hoch, oben
  \item[Hyperglykämie] hoher Blutzucker
  \item[Hypertonie] Bluthochdruck
  \item[Hypo-] tief, unten
  \item[Hypoglykämie] niedriger Blutzucker
  \item[Hypotonie] niedriger Blutdruck
  \item[Brady-] langsam
  \item[Bradykardie] langsamer Herzschlag
  \item[Bradypnoe] langsame Atmung
  \item[Tachy-] schnell
  \item[Tachykardie] schneller Herzschlag
  \item[Tachypnoe] schnelle Atmung
\end{description}
\end{multicols}
}{}{}{}{}

\Text{Beschreibung von Krankheiten und Symptomen}
{\begin{multicols}{2}
\begin{description}
    \item[akut] plötzlich; rascher Handlungsbedarf
    \item[chronisch] dauerhaft, wiederkehrend
    \item[\Z{rezidivierend}] \Z{wiederkommend, neuerlich}

\end{description}
\end{multicols}}
{}
{}
{}
{}


\Text{Stoffe}{
\begin{multicols}{2}
\begin{description}
  \item[Aminosäuren] Bausteine von Proteinen.
  \item[Base] \Lang{Syn.} \I{Lauge}. Stoff, welcher
  Protonen aufnehmen kann.\ZFootnote{Definition nach Brønstedt.}
  \Z{Sie reagieren mit Säuren und bilden Salze.}
  \cite{SchmuckChemie1}
  \item[Enzym] \label{topic:enzym}Hilfsstoff, welcher eine
  chemische Reaktion (Umwandlung von Stoffen in andere
  Stoffe) begünstigt.
  \item[Protein] Eiweiß. Proteine werden aus Aminosäuren
  zusammengesetzt.
  \item[Säure] Stoff, welcher Protonen abgeben kann. \Z{Sie
  reagieren mit Basen und bilden Salze.}
  \cite{SchmuckChemie1}.
  
\end{description}
\end{multicols}
}{}{}{}{}


\Text{Medizinische Fachrichtungen, Fachabteilungen
und spezielle Einrichtungen}{%
% \begin{multicols}{2}
\begin{description}
%   \item[\T{Anatomie}] Lehre vom Bau der
%   Körperteile~\cite{Pschy259}.
  \item[\T{Chirurgie}] Fachgebiet zur Erkennung und
  Behandlung von Erkrankungen, welche ohne operative Maßnahmen zu
  gesundheitlichen Schäden oder dem Tod führen würden
  \cite{Pschy259}. 
%   Diese sehr allgemeine Definition hat nur
%   mehr bedingte Gültigkeit, da viele der traditionell
%   chirurgischen Krankheitsbilder heutzutage konservativ
%   (ohne Operation) oder interventionell (kleiner,
%   minimal-invasiver Eingriff, ohne Eröffnung einer
%   Körperhöhle) behandelt werden können. 
  Man kann weitere
  Spezialgebiete unterscheiden, z.\,B.  Allgemeinchirurgie
  (Bauch), Gefäßchirurgie,
  Herzchirurgie, Thoraxchirurgie,
  Kinderchirurgie, Plastische Chirurgie etc.
  Daneben gibt es noch viele andere operative Fächer.
%   z.\,B. die Unfallchirurgie, Neurochirurgie, Gynäkologie,
%   Hals-Nasen-Ohrenheilkunde, Urologie etc.
\item[\T{Geburtshilfe}]
  \item[\T{Gynäkologie}] Frauenheilkunde
%   \item[\T{Hygiene}] Lehre von der Verhütung der Krankheiten
%   und der Erhaltung, Förderung und Festigung der Gesundheit
\item[\T{HNO}]
  \item[\T{ICU}] \Z{\Lang{engl.} \I{Intensive Care Unit};
  \I{Intensivstation}}. 
  Besondere Behandlungseinrichtung für intensive, oft lebenserhaltende
  medizinische Maßnahmen und Überwachungen, z.\,B. Beatmung,
  Dialyse, Sedierung, Therapie mit kreislaufwirksamen
  Substanzen etc. In der Regel unterscheidet man zwischen
  internistischen Intensivstationen, welche sich eher der
  Behandlung von Herz-Kreislauferkrankungen widmen, und
  allgemeinen oder anästhesiologischen Intensivstationen,
  mit meist eher chirurgischen Patienten.
  \item[\T{Innere Medizin}] \Lang{Syn.} \I{Interne Medizin}.
  Die Innere Medizin befasst sich mit der Vorbeugung,
  Diagnostik, konservativen und interventionellen Behandlung
  sowie Rehabilitation und Nachsorge von
  Gesundheitsstörungen und Krankheiten der inneren Organe
  und Organsysteme. Innerhalb der Inneren Medizin gibt es
  verschiedene Spezialgebiete: 
  \II{Kardiologie} (Herz), 
  \I{Hämatologie}  und \II{Onkologie} (Blut, Krebserkrankungen),
  \Z{\I{Angiologie} (Gefäße),
  \I{Endokrinologie} und \I{Diabetologie} (Stoffwechsel), 
  \I{Gastroenterologie} (Magen-Darm-Trakt), 
  \I{Nephrologie} (Niere),
  \I{Pneumologie} (\Lang{Syn.} \I{Pulmologie}, Lunge, eigenes Sonderfach), 
  \I{Rheumatologie} (entzündliche Erkrankungen des Bewegungsapparates)}. 
  
  \item[\T{Neurologie}] Fachgebiet der Medizin, welches sich 
  mit der Erforschung, Diagnostik und Behandlung von
  Erkrankungen des \E{Nervensystems} und der Muskulatur
  befasst.
  \item[\T{Nuklearmedizin}] Fachgebiet, welches die
  Anwendung radioaktiver Substanzen und offener Radionuklide
  zur Diagnostik und Therapie umfasst.
  \item[\T{Onkologie}] Fachgebiet, welches sich mit der
  Erkennung, Diagnostik, Therapie und Nachsorge von
  bösartigen Tumorerkrankungen (\enquote*{Krebs}) befasst.
  Die Onkologie an sich ist ein Teilgebiet der Inneren
  Medizin, meist ist jedoch eine fachübergreifende
  Behandlung notwendig.
  \item[]
  \item[\T{Pädiatrie}] \I{Kinderheilkunde}, zuständig für
  Patienten bis zur Vollendung des 18. Lebensjahres.
  \item[PCI] \Lang{Abkz.} \emph{\Z{Perkutane
  Coronar-Intervention}}; Herzkatheterbehandlung, siehe PTCA
  \item[\T{Physikalische Medizin}]
  \item[\T{Psychiatrie}] Die Psychiatrie beschäftigt sich
  mit der Prävention, Diagnostik und Therapie psychischer Störungen.
  In der
  Psychiatrie ist es besonders schwierig, zwischen
  \enquote*{gesund} und \enquote*{krank} treffsicher zu
  unterscheiden. Eine andere, tendenziell humorvolle
  Definition besagt daher:
  \Speech{\enquote{Der, der den Schlüssel
  hat, ist der Arzt.}  (Otto Pötzl \cite{Jolesch2})}
  \item[PTCA] \Lang{Abkz.} \emph{\Z{Perkutane Transluminale
  Coronar-Angiographie}}; Herzkatheteruntersuchung
  \item[\T{Pulmologie}] \Z{\Lang{Syn.}
  Pneumologie}. Lungenheilkunde
  
  \item[\T{Radiologie}] Fachgebiet, welches die
  \I{Bildgebung} mittels verschiedener Verfahren ermöglicht.
  Die Bildgebung kann auch zu therapeutischen Zwecken genutzt
  werden (interventionelle Radiologie).
  \item[\T{Stroke Unit}] Spezialstation zur Behandlung von
  \E{Schlaganfällen}
  \item[\T{Toxikologie}] Lehre der Vergiftungen und
  Giftstoffe
  \item[\T{Unfallchirugie} und \T{Traumatologie}] Lehre und
  Fachgebiet der Medizin, welches sich mit Unfällen und Verletzungen beschäftigt.
  \item[\T{Urologie}] Heilkunde des harnableitenden Systems
\end{description}
% \end{multicols}
}{}{}{}{}


\Text{Sonstige Begriffe}{
% \begin{multicols}{2}
\begin{description}
  \item[Bolus] \begin{inparaenum}
    \item \enquote{Happen}; 
    \item Stoßweise
  Verabreichung eines Medikaments (intravenös)
  \end{inparaenum}
  \item[Cave] \enquote{Achtung!}, \enquote{Hüte Dich!}
  \item[Obstruktion] Verlegung, Verengung. z.\,B. im Rahmen
  der \Abk{COPD }(\prettyref{topic:copd}).
  \opt{STATUS-EXPERIMENTAL}{\item[Organ] \TODO{GABS}{yyyy-mm-dd}{EIN /
  Begriffe: \enquote*{Organ} muss noch erklärt werden.}{}}
  \item[Pathologie, pathologisch, Patho-] Lehre von den
  Krankheiten. Als pathologisch wird alles Krankhafte
  bezeichnet. Die Vorsilbe \enquote*{Patho} weist auf etwas
  Abnormales, Krankhaftes hin.
  \item[Zyanose] Blaufärbung der Haut aufgrund eines
  mangelnden Sauerstoffgehaltes des Blutes. Betroffen sind
  insbesonders die Finger, Lippen und das Gesicht.
\end{description}
% \end{multicols}
}{}{}{}{}


\Text{Richtungs- und Lageangaben}{
Für die Beschreibungen von Richtungen und Lagebeziehungen
wird eine systematische Terminologie verwendet. Die
einschlägigen Bezeichnungen sind auch regelmäßig in Arzt-
und Entlassungsbriefen anzutreffen.
\begin{multicols}{2}
\Z{%
\begin{description}
  \item[lateral] \Lang{Abk.} lat.; seitlich
  \index{lateral|TLike}
  \index{lat.|see{lateral}}
  \item[medial] \Lang{Abk.} med.; mittig
  \index{medial|TLike}
  \index{med.|see{medial}}
  \item[dorsal] rückenwärts
  \index{dorsal|TLike}
  \item[ventral] bauchwärts
  \index{ventral|TLike}
  \item[cranial] kopfwärts
  \index{cranial|TLike}
  \item[caudal] fußwärts \Z{(eigentlich: schwanzwärts)}
  \index{caudal|TLike}
  \item[distal] \Lang{Abk.} dist.; körperfern
  \index{distal|TLike}
  \index{dist.|see{distal}}
  \item[proximal] \Lang{Abk.} prox.; körpernah
  \index{proximal|TLike}
  \index{prox.|see{proximal}}
  \item[dexter, -tra, trum] \Lang{Abk.} dext.; rechts
  \index{dexter|TLike}
  \index{dext.|see{dexter}}
  \item[sinister, -tra, trum] \Lang{Abk.} sin.; links
  \index{sinister|TLike}
  \index{sin.|see{sinister}}
\end{description}
}

\end{multicols}
}{}{}{}{}


\begin{Danger}
  \item Seiten- und Richtungsangaben (rechts, links, \ldots)
  beziehen sich immer auf die betreffende Person (Patienten)!
\end{Danger}

\subsection{\HeadingKeyword{Diagnose}} 
\index{Diagnose}
\label{topic:diagnose}
%Der Begriff Diagnose wird dauernd verwendet, dennoch ist er
%nicht einfach zu erklären. 
Eine Diagnose gibt einem Zustand
eine Bezeichnung, das heißt einem Beschwerdebild wird ein
definiertes Krankheitsbild zugeordnet.

\begin{center}\TableBaseModification
\noindent\begin{tabularx}{\linewidth}{XcX}\toprule
Beschwerden und Symptome des Patienten
&
{\color{darkBlue}\sffamily\bfseries $\rightarrow$~DIAGNOSE~$\leftarrow$}
&
Definiertes Krankheitsbild
\\\addlinespace
Patient
&
{\color{darkBlue}\sffamily\bfseries $\rightarrow$~DIAGNOSE~$\leftarrow$}
&
Lehrbuch
\\\addlinespace
Realität
&
{\color{darkBlue}\sffamily\bfseries $\rightarrow$~DIAGNOSE~$\leftarrow$}
&
Lehre
\\\bottomrule
\end{tabularx}
\end{center}

\begin{Synopsis}
  \SynopsisItem{Die Diagnose ordnet die Beschwerden des Patienten einem definierten
  Krankheitsbild zu.}
\end{Synopsis}

\TextSlide{Arten von Diagnosen}
{Diese Zuordnung kann einfach sein, sehr oft ist die Diagnosestellung
aber unsicher. Um die Unsicherheit dieser Diagnosen
anzuzeigen haben sich verschiedene Begriffe eingebürgert: 

\begin{itemize}
    \item \T{Verdachtsdiagnose}: 
    Die Diagnose ist unsicher, man hat zwar den Verdacht in
    Richtung  eines bestimmten Krankheitsbildes, andere
    ähnliche Krankheitsbilder können  allerdings nicht mit
    hinreichender Sicherheit ausgeschlossen werden. Die
    endgültige Klärung überlässt man jemand
    anderen\footnote{z.\,B. dem Krankenhaus}. \par Sanitäter
    erstellen i.\,d.\,R. nur Verdachtsdiagnosen.
     
    \item \T{Arbeitsdiagnose}: 
    Arbeitsdiagnosen sind Verdachtsdiagnosen, bei denen eine
    endgültige Klärung nicht  möglich ist. Für die
    Behandlung wird die wahrscheinlichste oder gefährlichste
    Diagnose ausgewählt um \Speech{\enquote{am Patienten
    arbeiten zu können}.}
    \item \T{Notfalldiagnose}: Eine
    Notfalldiagnose ist im Prinzip eine Arbeitsdiagnose. Der Patient ist dabei 
    in einem Zustand bei dem \EE{sofortiges Handeln} notwendig ist und die 
    eigentliche Diagnose, die zu dieser Situation geführt hat, in dem Moment 
    unwichtig ist. 
    \par Ein Beispiel hierfür sind die Bewusstlosigkeit und der
    Kreislaufstillstand. Die Ursachen dafür sind erstmal unwichtig, mit der
    entsprechenden Behandlung der Notfalldiagnose muss sofort begonnen werden.
    \item \T{Status post},  \T{Zustand nach}:
    \index{St.\,p.|see{Status post}}
    \index{Z.\,n.|see{Zustand nach}} \Lang{Abk.}
    \Speech{\enquote*{St.\,p.}}, \Speech{\enquote*{Z.\,n.}}. Bezeichnet eine
    \enquote*{alte Diagnose}, d.\,h. eine Erkrankung oder eine Behandlung
    die einmal durchgemacht wurde. Z.\,B.:
    \begin{itemize}
        \item \Tt{Z.\,n. Blinddarm-Operation}: Der Patient
        hatte einmal eine Blinddarm-Operation.
        \item \Tt{Z.\,n. Herzinfarkt 2/2006}: Der Patient
        hatte im Februar 2006 einen Herzinfarkt.
    \end{itemize}
    Umgangssprachlich wird oft auch der Umstand der zu einer
    Diagnose geführt hat mit \Speech{\enquote*{Z.\,n.}}
    angegeben.
    \begin{itemize}
        \item \Tt{V.\,a. Schenkelhalsfraktur re., Z.\,n. Sturz}: Der
        Patient ist gestürzt und hat deshalb möglicherweise eine
        Schenkelhalsfraktur rechts.
    \end{itemize}
\end{itemize}}
{
\begin{itemize}
    \item {Verdachtsdiagnose}
    \item {Arbeitsdiagnose}
    \item {Notfalldiagnose}
    \item {Status post}, {Zustand nach}
\end{itemize}}
{}
{}
{}




\TextSlide{Diagnosecodes}
{Zu statistischen und organisatorischen Zwecken werden
Diagnosen häufig zur weiteren Verarbeitung und Auswertung
codiert, d.\,h. die jeweilige Diagnose wird einem
definierten Code zugeordnet. 

Ein weit verbreitetes System
ist ist die \I{Internationale statistische Klassifikation
der Krankheiten und verwandter Gesundheitsprobleme}
(\T{ICD}\footnote{\Lang{engl.} International Statistical
Classification of Diseases and Related Health Problems}),
welche von der Weltgesundheitsorganisation (WHO)
herausgegeben wird. Derzeit aktuell ist die Version 10
({\T{ICD-10}}). \Z{Ein ICD-10-Code beginnt mit einem
Buchstaben. Die ersten drei Stellen geben eine grobe
Diagnose an. Stellen 4 und 5 werden von den ersten drei
Stellen mit einem Punkt getrennt und erlauben eine genauere
Beschreibung der Diagnose.}

\Z{\EE{Beispiel}:

\begin{itemize}
  \item \EE{\texttt{S50-S59}}: Verletzungen des Ellenbogens
  und des Unterarmes
  \item \EE{\texttt{S52.3-}}: Fraktur des Radiusschaftes
  \item \EE{\texttt{S52.31}}: Fraktur des distalen
  Radiusschaftes mit Luxation des Ulnakopfes
\end{itemize}}}
{\begin{itemize}
  \item ICD-10
\end{itemize}}
{}
{}
{}

\Text{Differentialdiagnosen (DD)}
{\index{Differentialdiagnose}%
Diagnosen, welche bei den jeweiligen Symptomen auch
möglich, jedoch weniger wahrscheinlich sind wie jene
Diagnose, für welche man sich \enquote{entscheidet}.

\begin{quote}
Beispiel: Patient mit Kopfschmerzen nach ausgiebigen
Alkoholgenuss am Vortag.
\begin{description}
    \item[Verdachtsdiagnose:] \E{\enquote{Kopfschmerzen}}
    \item[Differentialdiagnosen:]
    Migraine, Hirnblutung, Schlaganfall,
    Hypoglykämie, Intoxikation, \ldots
\end{description}
\end{quote}}
{}
{}
{}
{}






\bigskip

% \clearpage % TEMP temp clearpage
\section{Wichtige \HeadingKeyword{Pathomechanismen}}


%\subsection{Das \HeadingKeyword{Gewebe} betreffende Störungen}
\Text{{\TxBeschreibung}}
{%
Durch krankhafte Vorgänge kann es im Gewebe zu
\begin{itemize}
    \item \T[Oedem@Ödem]{Ödemen} (Wasseransammlung im
    Gewebe),
    \item Emphysemen (Luftansammlung im Gewebe),
    \item Entzündungen (\FullPrettyRef{topic:entzuendung}),
    \item Tumoren (\FullPrettyRef{topic:tumor}) oder
    \item Nekrosen (Absterben von Zellen bzw. Gewebe)
\end{itemize}
kommen. Die letzten drei Punkte werden im Folgenden näher
erklärt.
}
%{%
%
%}



\TextSlide
{\HeadingKeyword{Entzündung} ist eine Reaktion des Körpers}
{%
%%%
\index{Entzündung}%
\label{topic:entzuendung}%
%%%
\Z{\protect\Lang{engl., lat.} inflammation.}
\DefinitionInPara{Eine \T{Entzündung} ist eine 
Reaktion des Körpers welche charakterisiert
ist durch \begin{inparaenum}
    \item \EE{Rötung},
    \item \EE{Schwellung},
    \item \EE{Überwärmung},
    \item \EE{Schmerz} und
    \item \EE{Funktionsverlust}.
\end{inparaenum}}
Mögliche Auslöser sind Fremdkörper, Krankheitserreger,
mechanische Beanspruchung, chemische Substanzen oder andere
pyhsikalische Noxen.
Typisch für die Bezeichnung von Entzündungen ist das
\EE{Suffix {\enquote*{\T{-itis}}}}\index{itis@-itis},
welches i.\,d.\,R. an die betreffende Organbezeichnung angehängt
wird: 
\begin{compactitem}
  \item Gastritis ${\rightarrow}$ \enquote*{Magenentzündung}
  \item Hepatitis ${\rightarrow}$ Leberentzündung
  \item Pankreatitis ${\rightarrow}$ Entzündung der Bauchspeicheldrüse,
  \item Meningitis ${\rightarrow}$ Hirnhautentzündung
  \item Rhinitis ${\rightarrow}$ \enquote*{Schnupfen}
  \item Otitis ${\rightarrow}$ Ohrenentzündung
  \item \ldots
\end{compactitem}
}
{
\begin{itemize}
  \item Reaktion des Körpers
  \begin{enumerate}
    \item \EE{Rötung},
    \item \EE{Schwellung},
    \item \EE{Überwärmung},
    \item \EE{Schmerz} und
    \item \EE{Funktionsverlust}.
  \end{enumerate}
  \item Auslöser
  \begin{itemize}
    \item Fremdkörper
    \item Krankheitserreger
    \item Mechanische Beanspruchung
    \item Chemische Substanzen
    \item Physikalische Noxen
    \item \ldots{}
  \end{itemize}
  \item {\enquote*{{-itis}}}
\end{itemize}


}
{}
{}
{}


\Text{Ischämie}
{%
\label{topic:ischaemie}%
\DefinitionInPara{\T{Ischämie}  bezeichnet die akute
Unterversorgung eines Gewebes mit sauerstoffreichem Blut infolge der Verstopfung
des dafür zuständigen Blutgefäßes.} Eine Ischämie kann
grundsätzlich überall im Körper vorkommen. Es kommt zu
Funktionsausfällen des Gewebes bishin zum Gewebstod.
Schmerzen sind ein typisches Symptom. Häufig betroffene
Gewebe sind:

\begin{itemize}
  \item Herzmuskel (Herzinfarkt,
  \FullPrettyRef{topic:herzinfarkt})
  \item Hirn (trockener Insult,
  \FullPrettyRef{topic:insult})
  \item Darm (Mesenterialinfarkt,
  \FullPrettyRef{topic:mesenterialinfarkt})
  \item Extremitäten (Periphere arterielle
  Verschlusskrankheit, arterieller peripherer
  Gefäßverschluss,
  \FullPrettyRef{topic:arterieller-gefaessverschluss})
\end{itemize}
}
%{%
%
%}



\TextSlide{Tumore, Neubildungen, \HeadingKeyword{Krebs}}
{%
\label{topic:tumor}%
Die meisten Zellen des menschlichen Körper teilen sich
regelmäßig. Diese Form der Zellvermehrung ist für das
Überleben des Menschen notwendig. Eine Kontrolle dieses
Vorgangs ist jedoch erforderlich um zu gewährleisten dass
die Körperzellen nicht wild zu wuchern beginnen. Der Körper
verfügt über verschiedene Mechanismen um die Zellteilung im
richtigen Maß zu halten. 

Wenn diese Mechanismen versagen, kommt es zu einer
\E{Entartung} und \E{Wucherung}, meist in Form einer
\EE{Raumforderung} und Schwellung (\T{Tumor}\footnote{Der
Begriff \E{\enquote*{Tumor}} bedeutet ganz allgemein
\enquote*{Schwellung}. Umgangssprachlich wird er oft
(fälschlicherweise) mit \enquote*{Krebs} gleichgesetzt.}).
Man spricht dann von einer Neubildung \Z{(Neoplasie)}, in
Arztbriefen oft mit \enquote{{\ttfamily N.}} abgekürzt. Der
medizinische Fachbereich der sich mit Neubildungen
(Krebserkrankungen) befasst ist die \E{Onkologie}.
Grundsätzlich können alle teilungsfähigen Zellen entarten.
Dies kann spontan geschehen oder durch andere Faktoren
(radioaktive Strahlung, chemische Substanzen) begünstigt
bzw. ausgelöst werden.

\Z{Man unterscheidet zwischen gutartigen \Z{(benignen)} und
bösartigen \Z{(malignen)} Neubildungen. Gutartige
Neubildungen \enquote*{bleiben wo sie sind} und wachsen vor
sich hin. Bösartige Neubildungen bilden Absiedelungen, sog.
\E{Metastasen}, welche z.\,B. über die Blutbahn an andere
Stellen des Körper gelangen und dort zu wuchern beginnen.
Je nach Zellart unterscheidet man bei bösartigen
Wucherungen zwischen \I[Karzinom]{Karzinomen},
\I[Sarkom]{Sarkomen} und \I[Blastom]{Blastomen}, bei
Krebserkrankungen der Blutzellen spricht man von
\I[Leukämie]{Leukämien} bzw. \I[Lymphom]{Lymphomen}.}

\begin{table}[H]\TableBaseModification
\Captionof{table}{Beispiele: Häufige Neubildungen.}

\begin{tabularx}{\linewidth}{XXl}
\toprule\addlinespace
\EE{Schlagwort}
&
\EE{Beispiel}
&
\TabHeading{\Z{Im Arztbrief}}
\\\addlinespace\midrule\addlinespace
Lungenkrebs
&
Bronchialkarzinom
&
{\ttfamily\Z{N. bronchii}}
\\\addlinespace
Brustkrebs
&
Mammakarzinom
&
{\ttfamily\Z{N. mammae}}
\\\addlinespace
Darmkrebs
&
Rektumkarzinom
&
{\ttfamily\Z{N. recti}}
\\\addlinespace
Blutkrebs
&
Leukämie, versch. Arten
&
{\ttfamily\Z{ALL, AML, CML, \ldots}}
\\\addlinespace\bottomrule
\end{tabularx}
\end{table}

\subparagraph{Was ist das Problem?}

Zusammengefasst ergeben sich drei wesentliche Probleme:
\begin{enumerate}
    \item \EE{Neubildungen brauchen Raum.} Dieser Raum ist
    nicht endlos vorhanden. Je nach Lage können wichtige
    Strukturen beengt oder behindert werden (Hirn,
    Luftröhre, \ldots).
    \item \EE{Funktionsverlust des gesunden Gewebes.}
    Durch den Befall mit (nutzlosen) Metastasen wird das
    gesunde Gewebe soweit durchwachsen und geschwächt, dass
    es seine Funktion nicht mehr richtig ausführen kann
    (z.\,B. Knochenbefall).
    \item \EE{Schmerz.} Der Befall verursacht oft starke
    Schmerzen. Krebspatienten benötigen oft eine
    dauerhafte, starke Schmerztherapie.
\end{enumerate}

\subparagraph{Therapie}

Der Therapieerfolg ist sehr stark von der genauen Art der
Neubildung und vom jeweiligen Patienten abhängig und reicht
von sehr gut bis nicht vorhanden. 
Im wesentlichen kommen bei der Krebstherapie folgende
Methoden zum Einsatz:
\begin{itemize}
    \item \E{Chemotherapie}: Einsatz von Medikamenten welche
    die Zellteilung unterdrücken bzw. Zellen absterben
    lassen.
    \item \E{Bestrahlung}: zur Abtötung oder Vorbeugung von
    Geschwülsten.
    \item \E{Operation}: Chirurgische Entfernung der
    Neubildung(en).
\end{itemize}
Diese Therapien sind für den Körper sehr anstrengend und
haben viele Nebenwirkungen, wie z.\,B. Haarausfall, Übelkeit,
Erbrechen, \EE{Schwächung des Immunsystems}, Nervenschäden
u.\,v.\,a.\,m.

\cite{UlfigHistologie1,GantenTumorerkr98,Uicc3}
}
{%

}





\ChapterEnd

